%!TeX program = xelatex
\documentclass{SYSUReport}
\usepackage{tabularx} % 在导言区添加此行
\usepackage{float}
% 根据个人情况修改
\headl{}
\headc{}
\headr{并行程序设计与算法实验}
\lessonTitle{并行程序设计与算法实验}
\reportTitle{Lab6-Pthreads并行构造}
\stuname{xxx}
\stuid{xxx}
\inst{计算机学院}
\major{xxx}
\date{2025年4月30日}

\begin{document}

% =============================================
% Part 1: 封面
% =============================================
\cover
\thispagestyle{empty} % 首页不显示页码
\clearpage

% % =============================================
% % Part 4： 正文内容
% % =============================================
% % 重置页码，并使用阿拉伯数字
% \pagenumbering{arabic}
% \setcounter{page}{1}

%%可选择这里也放一个标题
%\begin{center}
%    \title{ \Huge \textbf{{标题}}}
%\end{center}

\section{实验目的}
\begin{itemize}
  \item 深入理解C/C++程序编译、链接与构建过程
\item 提高Pthreads多线程编程综合能力
\item 并行库开发与性能验证
\end{itemize}

\section{实验内容}
\begin{itemize}
 \item 基于Pthreads的并行计算库实现：parallel\_for+调度策略
\item 动态链接库生成和使用
\item  典型问题的并行化改造：矩阵乘法，热传导问题
\end{itemize}

\section{实验过程}

\subsection{环境与工具}
简要说明实验所使用的操作系统、编译器 (gcc/g++) 版本、以及 Pthreads 库。
\begin{itemize}
    \item 回答：
\end{itemize}
\subsection{核心函数实现 (\texttt{parallel\_for})}
简要描述 \texttt{parallel\_for} 函数的关键设计思路和实现。重点说明线程创建、任务划分和同步机制。
\begin{lstlisting}[language=C++, caption={parallel\_for 函数核心代码片段}]
// 在此处粘贴你的 parallel_for 函数核心代码
// 简单注释说明关键部分即可
\end{lstlisting}
\subsection{动态库生成与使用}
说明生成动态链接库 (.so) 的主要命令或 Makefile 规则，并简述如何在主程序 (如矩阵乘法、热传导) 中链接和调用该库。
\begin{verbatim}
# 示例：生成动态库命令
...
# 示例：编译链接主程序
...
\end{verbatim}
\subsection{应用测试 (热传导)}
简述如何将 \texttt{parallel\_for} 应用于热传导问题，替换原有的并行机制。描述测试设置，如网格大小和线程数。
\begin{itemize}
    \item 回答：
\end{itemize}

\section{实验结果与分析}

\subsection{性能测试结果}
展示不同线程数和调度方式下，自定义 Pthreads 实现与原始 OpenMP 实现的性能对比。
\begin{table}[h!]
\centering
\caption{热传导问题性能对比 (Pthreads vs OpenMP, 网格大小: M x N)}
% Updated column definition: 4 columns total
\begin{tabular}{|c|c|c|c|}
\hline
\textbf{线程数} & \textbf{调度方式} & \textbf{自定义 Pthreads} & \textbf{原始 OpenMP} \\
\cline{3-4} % Line under Pthreads and OpenMP Time columns
 & \textbf{(Pthreads)} & 时间 (s) & 时间 (s) \\
\hline
1 (串行) & N/A & T1\_pth & T1\_omp \\
\hline
% Updated colspan to 4
\multicolumn{4}{|c|}{Pthreads: 静态调度 (Static)} \\
\hline
% Removed Speedup columns
2 & Static & T2\_pth\_static & T2\_omp \\
4 & Static & T4\_pth\_static & T4\_omp \\
8 & Static & T8\_pth\_static & T8\_omp \\
16 & Static & T16\_pth\_static & T16\_omp \\
... & ... & ... & ... \\
\hline
% Updated colspan to 4
\multicolumn{4}{|c|}{Pthreads: 动态调度 (Dynamic, chunk=X)} \\
\hline
% Removed Speedup columns
2 & Dynamic & T2\_pth\_dynamic & T2\_omp \\
4 & Dynamic & T4\_pth\_dynamic & T4\_omp \\
8 & Dynamic & T8\_pth\_dynamic & T8\_omp \\
16 & Dynamic & T16\_pth\_dynamic & T16\_omp \\
... & ... & ... & ... \\
\hline
\end{tabular}
\label{tab:performance_comparison_time_only}
\end{table}

\subsection{结果分析与总结}
简要分析性能结果，说明是否达到了并行加速效果。
\begin{itemize}
    \item 回答：
\end{itemize}
注：实验报告格式参考本模板，可在此基础上进行修改；实验代码以zip格式另提交；最终提交内容包括实验报告(pdf格式)和实验代码(zip压缩包格式)
\end{document}